\subsection{Segunda rodada: parecer sobre todas as indicações recebidas}
A segunda rodada retoma individualmente as 66 indicações preservadas no catálogo original, sem reduzir esse universo às fontes que deram resultados positivos. O parecer completo está no apêndice e em \texttt{data/references\_audit.csv}. A consulta combinou busca por título e período, busca literal de âncoras das 22 frases, aquisição de exemplares acessíveis, pesquisa de OCR e conferência visual de passagens selecionadas. Não houve leitura visual integral de todas as páginas de todos os títulos. A busca de uma frase sem resultado e a impossibilidade de baixar um documento não constituem demonstração de inexistência.

Os 66 pareceres distribuem-se em cinco estados mutuamente exclusivos: três indicações com representação do riso autenticada; sete com referência lexical ou temática identificada; dezenove com identificação bibliográfica, sem autenticação de passagem; quatro com exemplar específico recuperado, mas sem autenticação da frase atribuída; e trinta e três sem ocorrência autenticada ou com identificação insuficiente. São contagens do andamento da auditoria, não estimativas de prevalência histórica. A categoria lexical inclui uma verificação contextual de \textit{Policarpo Quaresma}, discutida adiante; não significa que a fala atribuída seja expressão de riso.

As três indicações positivas são \textit{Urupês}, \textit{Diário Nacional} e \textit{Macunaíma}. Em \textit{Urupês}, quinta edição, de 1919, o conto \textit{O engraçado arrependido}, página impressa 19, correspondente à página 24 do PDF, traz \forma{Quá, quá, quá!}, seguida de descrição explícita dos risos do grupo \cite{lobato1919}. O exemplar confirma a grafia em 1919; a atribuição da mesma passagem à primeira edição de 1918 exigiria cotejo específico. A descoberta amplia a distribuição documental de \forma{quá} anterior a 1935, sem demonstrar filiação a \forma{kkkk}.

No \textit{Diário Nacional}, número 290, de 17 de junho de 1928, página 4, o discurso de Mattos Pimenta contém \forma{(risos, muito bem)} \cite{diarionacional1928}. A data foi conferida no cabeçalho e a rubrica na imagem da página. O dado confirma uma representação descritiva de reação da audiência, que deve permanecer distinta da imitação de sua sonoridade. Sua existência não autentica qualquer sequência \forma{ah} atribuída genericamente aos jornais da década.

Na edição de \textit{Macunaíma} publicada pela Editora UFFS em 2019, a página 139 contém \forma{Quá! quá! quá quaquá!...}, com indicação imediata de uma grande gargalhada \cite{andrade2019}. A ficha catalográfica, no final do PDF, identifica a edição consultada. A primeira publicação da obra em 1928 e a data do testemunho efetivamente examinado são registradas separadamente: este item entra na cobertura anual em 2019, sem retrodatação automática da passagem. Trata-se de ocorrência diferente da fala com \forma{Ah! ah! ah!} apresentada no roteiro.

A pesquisa também recuperou um documento adicional: \textit{Jornal de Notícias}, de São Paulo, número 267, de 2 de março de 1947, página 1. Na notícia sobre a fala de Bidault, encontra-se \forma{(Risos dos presentes)} \cite{jornalnoticias1947}. O documento surgiu em uma busca por \textit{Folha da Manhã}; a inspeção do título demonstrou que era outro jornal. Por isso, foi incluído como testemunho autônomo e não como confirmação da indicação sobre a \textit{Folha}. O inventário passa de quinze para dezenove atestações: dez com imagem conferida, seis com transcrição primária eletrônica e três interações reportadas em pesquisas.

\subsubsection{Referências ao riso que não autenticam uma grafia}
Em \textit{Canaã}, o PDF consultado traz uma referência humana ao riso na página 28 \cite{r2canaa}. Na segunda edição de \textit{Recordações do escrivão Isaías Caminha}, de 1917, a página 26 do arquivo contém uma forma de \forma{rir} seguida da menção a uma pilhéria \cite{r2isaias}. Em \textit{Esfinge}, capítulo VI, a transcrição consultada apresenta risos descritos na cena \cite{r2esfinge}. Em \textit{Eu}, o poema \textit{Monólogo de uma sombra} associa choro e riso \cite{r2eu}. São ocorrências lexicais ou temáticas. Elas não se convertem em prova de uma sequência onomatopaica específica.

No exemplar de \textit{Memórias sentimentais de João Miramar}, de 1924, página impressa 114, PDF 116, encontra-se uma referência ao desaparecimento do sorriso dos presentes \cite{r2miramar}. Nas provas da segunda edição de \textit{São Bernardo}, de 1938, a página 63 do PDF contém a expressão \forma{sorrisos frios} \cite{r2saobernardo}. A busca no texto extraído não localizou a âncora \forma{latim de missa}; a qualidade do reconhecimento de caracteres e a configuração das provas impedem transformar o resultado em prova de ausência. O exame tampouco permite datar automaticamente essa referência lexical de 1934.

O cotejo de \textit{Triste fim de Policarpo Quaresma} apresenta resultado mais específico. Na terceira parte, capítulo II, a fala de Floriano é \forma{Você, Quaresma, é um visionario...}, na grafia do texto eletrônico consultado \cite{r2policarpo}. A passagem não contém a risada nem a continuação alegadas no roteiro, e a ordem dos termos difere da transcrição recebida. Há, portanto, um núcleo reconhecível de fala literária, mas isso não autentica a citação integral apresentada como documento do riso. A versão integral dessa citação continua excluída do corpus de atestações.

\subsubsection{Identificação bibliográfica e correções de cronologia}
Catálogos e páginas editoriais cumprem uma função distinta da leitura de uma passagem. A BBM-USP identifica \textit{A menina morta} em edição de 1954 \cite{r2meninamorta}. A alegação de que a obra teria sido escrita nos anos 1920 não foi demonstrada. A bibliografia da Academia Brasileira de Letras distingue a edição italiana de \textit{Zero}, de 1974, e a edição em português, de 1975 \cite{r2zero}; mesmo que fosse comprovado trabalho preparatório nos anos 1960, esse fato não equivaleria à circulação de uma grafia naquele período.

A fonte institucional da Universidade do Estado do Pará situa a primeira publicação de \textit{Chove nos campos de Cachoeira} em 1941 \cite{r2chove}. O ano de 1998 informado no roteiro pode corresponder a edição posterior, que precisaria ser individualizada. A página editorial de \textit{Budapeste} registra lançamento em 10 de setembro de 2003, enquanto a de \textit{Estorvo} registra 2 de agosto de 1991 \cite{r2budapeste,r2estorvo}. Suposto planejamento de \textit{Budapeste} no final dos anos 1990 não foi documentado e não pode servir como data de uma ocorrência impressa.

Foram localizadas identificações bibliográficas de \textit{O quinze}, \textit{Fogo morto}, \textit{O continente}, \textit{Sagarana} e \textit{Grande sertão: veredas} \cite{r2quinze,r2fogo,r2continente,r2sagarana,r2grande}. O mesmo vale para \textit{Vestido de noiva}, \textit{Quarup}, \textit{A paixão segundo G.H.}, \textit{As meninas}, \textit{Sargento Getúlio} e \textit{A morte de D.J. em Paris} \cite{r2vestido,r2quarup,r2paixao,r2meninas,r2getulio,r2morte}. Foram também consultados registros editoriais de \textit{Viva o povo brasileiro}, \textit{O Alquimista} e \textit{Feliz ano velho} \cite{r2viva,r2alquimista,r2feliz}. Essas referências são citadas como instrumentos de identificação, com seu alcance explicitado na bibliografia. Não substituem o texto integral nem confirmam as frases sobre riso que lhes foram atribuídas. No caso do teatro, escrita, estreia e edição em livro não são uma mesma data.

\subsubsection{Exemplares recuperados e limites das buscas negativas}
A edição de 25 de julho de 1904 da \textit{Gazeta de Notícias}, número 207, foi recuperada e teve o cabeçalho conferido \cite{r2gazeta1904}. O PDF do número 206, inicialmente candidato, é de 24 de julho; a diferença foi resolvida pela imagem, não pela hipótese construída a partir do endereço eletrônico. A busca nas quatro páginas do OCR do número 207 não autenticou a frase do ovo. Essa informação descreve o procedimento realizado; não demonstra que uma crônica semelhante jamais tenha sido publicada.

Também foram recuperados \textit{O Malho}, número 526, de 12 de outubro de 1912, e \textit{Diário de Notícias}, número 6564, de 18 de março de 1944, ambos com data conferida no cabeçalho \cite{r2malho1912,r2diario1944}. A pesquisa no OCR não autenticou, respectivamente, a frase do prefeito e da carne barata e a declaração ministerial sobre racionamento. A revisão visual integral das 58 e 12 páginas, sobretudo de desenhos e legendas, permanece necessária para um resultado negativo mais forte. O PDF identificado pelo acervo como número 23 de \textit{Fon-Fon!} de 1935 foi adquirido, mas o expediente e a frase atribuída à data de 8 de junho não foram autenticados.

Uma notícia de \textit{Fon-Fon!}, de maio de 1925, página 3, informa o retorno do \textit{Correio da Manhã} após suspensão policial \cite{r2fon1925}. O dado recomenda revisar a alegação de um trecho publicado em 14 de fevereiro daquele ano, mas não basta, sozinho, para delimitar todos os dias de suspensão. A frase da baronesa continua sem autenticação; não se declara inexistente a edição apenas por esse indício.

Algumas indicações não individualizam uma fonte: \textit{Reforma Agrária}, de ``Marquie'', não teve identificação bibliográfica suficiente; \textit{Os segredos do Orkut}, como livro de 2005, não foi identificado por autor e editora; \textit{Manual do MSN Messenger} não veio acompanhado de autoria ou ISBN; e \textit{Literatura de Blog} designa uma categoria. A recuperação de postagem homônima de 2008 não autentica um livro de 2005. A falha HTTP 404 de uma URL tentada para \textit{O Pasquim}, número 12, tampouco demonstra que o fascículo não exista. Bloqueios, limites de requisições e falhas de OCR foram preservados no manifesto.

O resultado final desta rodada é um parecer para cada indicação e uma ampliação documentada do corpus, mantendo os casos inconclusivos visíveis. Nenhuma das 22 frases originalmente apresentadas como transcrições foi autenticada integralmente. Isso não significa que todas sejam inexistentes: significa que a documentação reunida não autoriza publicá-las entre aspas como testemunhos primários. As ocorrências efetivamente encontradas são descritas com suas próprias grafias, datas e localizadores.
